%% ============================================================
%% STRESZCZENIE — OPIS FORMATKI
%% ============================================================
\chapter*{Streszczenie}
\addcontentsline{toc}{chapter}{Streszczenie}

Niniejszy dokument jest szablonem \LaTeX{} dla prac dyplomowych realizowanych na Wydziale Elektroniki, Telekomunikacji i Informatyki Politechniki Gdańskiej. Struktura pracy powinna być zgodna z wytycznymi określonymi w \href{\detokenize{https://cdn.files.pg.edu.pl/eti/Dziekanat/regulaminy/ZR%2045-2024.pdf}}{Zarządzeniu Rektora nr 45/2024} z 15 listopada 2024 r. Wytyczne te definiują ogólne zasady przygotowania prac dyplomowych oraz wymagania edytorskie. Niniejsza formatka odzwierciedla, precyzuje i w wybranych miejscach rozszerza te zasady (m.in. dotyczące numerowania tabel, rysunków, wzorów i cytowań). Ułatwia również pracę z odwołaniami, podpisami i spisem treści. Zawiera wstępnie skonfigurowany zestaw pakietów \LaTeX{} oraz praktyczne wytyczne ,,dla autora'', tj. co umieszczać w poszczególnych rozdziałach pracy, jak opisywać metody i wyniki oraz w jaki sposób przygotowywać rysunki i tabele.

Głównym plikiem projektu jest \texttt{main.tex}. Ustawienia formatowania zebrano w oddzielnym pliku klasy \texttt{\_formatting.cls}, który zawiera wszystkie wymagane elementy budujące strukturę pracy. Dodatkowe wskazówki dotyczące formatowania pracy dyplomowej opisano w osobnych rozdziałach formatki. Struktura dokumentu obejmuje cztery główne rozdziały umieszczone w katalogu \emph{chapters} oraz dodatek z przykładowym kodem w języku Python. Każdy rozdział znajduje się w osobnym pliku \texttt{.tex}. Bibliografia jest zarządzana przy użyciu pakietu \texttt{biblatex} (\texttt{biber}) zgodnie z wytycznymi normy PN-ISO 690:2012, co zapewnia numerację w nawiasach kwadratowych zgodnie z kolejnością cytowań. Stronę tytułową oraz oświadczenie studenta dołącza się jako pliki PDF wygenerowane w systemie MojaPG i wstawia do dokumentu poleceniem \verb|\includepdf| w pliku głównym.

Całość projektu zoptymalizowano pod kompilację w środowisku Overleaf z koniecznym użyciem silnika LuaLaTeX (dystrybucja TeX Live w wersji co najmniej 2022). Domyślną czcionką zastosowaną w formatce jest Arial w rozmiarze 10 pt, a jej pliki zostały dołączone w katalogu \texttt{fonts}, co umożliwia kompilację projektu w dowolnym środowisku. Takie rozwiązanie łączy automatyzację formatowania z elastycznością \LaTeX{} oraz zapewnia pełną zgodność ze standardami Politechniki Gdańskiej. Należy pamiętać, że dokument musi spełniać wymogi przepisów antyplagiatowych — wygenerowany plik PDF nie może być w zakresie tekstu plikiem graficznym ani szyfrowanym; nie może także wykorzystywać niestandardowego kodowania czcionek (np. poprzez podmianę kodów znaków) ani „składania” znaków diakrytycznych z kilku znaków w jeden glif. Podsumowując, celem szablonu jest zminimalizowanie pracy technicznej autora, tak aby mógł on skupić się na treści, przy jednoczesnym zachowaniu jednolitego stylu edytorskiego wszystkich prac powstających na Wydziale ETI. Wszelkie uwagi dotyczące szablonu prosimy kierować do autora formatki: dr. inż. Sebastiana Dziedziewicza, e-mail: \textit{sebastian.dziedziewicz@pg.edu.pl}. Poniżej zamieszczono przykładowe słowa kluczowe, które mogą zostać wykorzystane w pracy.

\section*{Słowa kluczowe:}
\noindent uczenie maszynowe, sterowanie PID, systemy wbudowane, układy analogowe, modulacja, systemy rozproszone
\section*{Dziedzina nauki i techniki, zgodnie z wymogami OECD:}
\noindent nauki inżynieryjne i techniczne;  elektrotechnika, elektronika, inżynieria informatyczna; robotyka i automatyka; telekomunikacja.
